\documentclass[../../../include/open-logic-section]{subfiles}

\begin{document}

\olfileid{sth}{ordinals}{wo}
\olsection{良序}

基本概念如下：

\begin{defn}
关系 $<$ \emph{良序} $A$，当且仅当它满足以下两个条件：
\begin{enumerate}
	\item $<$ 是连通的，即对所有的 $a, b \in A$，$a < b$，或 $a = b$，或 $b < a$；
	\item $A$ 的每个非空子集都有一个 $<$-极小元，即如果 $\emptyset \neq X \subseteq A$ 那么 $(\exists m \in
	X)(\forall z \in X)z \nless m$
\end{enumerate}
\end{defn}

不难看出，前文考虑的三个例子确实都是良序关系。

\begin{prob}
\Olref[sth][ordinals][idea]{sec} 给出了自然数上的三个序关系例子。请验证每一个都是良序。
\end{prob}

以下是关于良序的几个基本但极其重要的观察。

\begin{prop}\ollabel{wo:strictorder}
如果 $<$ 良序 $A$，那么 $A$ 的每个非空子集都有唯一的 $<$-最小元，并且 $<$ 是反自反、非对称和传递的。
\end{prop}

\begin{proof}
如果 $X$ 是 $A$ 的非空子集，它有一个 $<$-极小元 $m$，即 $(\forall z \in X)z \nless m$。由于 $<$ 是连通的，$(\forall z \in X)m \leq z$。所以 $m$ 是 $X$ 的 $<$-最小元。

对于反自反性，取定 $a \in A$；$\{a\}$ 的 $<$-最小元是 $a$，所以 $a \nless a$。对于传递性，如果 $a < b < c$，那么由于 $\{a, b, c\}$ 有一个 $<$-最小元，可得 $a < c$。非对称性则由反自反性和传递性得出。
\end{proof}

\begin{prop}\ollabel{propwoinduction}
如果 $<$ 良序 $A$，那么对任意公式 $\phi(x)$：
\[
	\text{如果 }(\forall a \in A)((\forall b < a)\phi(b) \lif
		\phi(a))\text{，那么 }(\forall a \in A)\phi(a).
\]
\end{prop}

\begin{proof}
我们将证明其逆否命题。假设 $\lnot(\forall a \in
A)\phi(a)$，即 $X = \Setabs{x \in A}{\lnot\phi(x)} \neq
\emptyset$。那么 $X$ 有一个 $<$-极小元 $a$。因此 $(\forall b < a)\phi(b)$ 但 $\lnot \phi(a)$。
\end{proof}

\noindent 这最后一个性质应当让人想起自然数上的强归纳原理，即：如果 $(\forall n \in \omega)((\forall m < n)\phi(m)
\lif \phi(n))$，那么 $(\forall n \in \omega)\phi(n)$。而这个性质使得良序成为一个非常\emph{稳健}的概念。\footnote{提醒：所有公式都可以带参数（除非另有明确说明）。}

\end{document}